% ============================================
% CONT-04
% Gerado de molambudos.md — build_miolo.py
% ============================================
\section*{\textbf{CONT-04} --- Você Começa a Sonhar}
\markboth{CONT-04 --- Você Começa a Sonhar}{}

\noindent \textit{O olho no escuro}

\subsection*{Noite 1}

\noindent Você sonhou com o olho.

\noindent Não foi um sonho comum. Foi uma \textit{invasão}. Você estava num lugar que não reconhecia --- podia ser sua casa, podia ser o Colônia, podia ser a vala. O lugar mudava enquanto você olhava. As paredes respiravam.

\noindent E no canto, o olho.

\noindent Amarelo. Brilhante. Flutuando sozinho na escuridão.

\noindent Você queria acordar. Não acordou. O olho se aproximou. Você viu os vasos sanguíneos dentro dele --- não eram vasos humanos. Eram finos demais. Ramificados demais. Pareciam raízes.

\noindent O olho falou.

\noindent Não com som. A voz dele era uma \textit{sensação}. Uma compressão no peito. Um aperto na garganta. A certeza de que alguém estava atrás de você, mesmo você sabendo que não havia ninguém.

\noindent \textit{\textquotedblleft{}Você está lendo.\textquotedblright{}}

\noindent Você tentou responder que sim. Não conseguiu. A boca não abria.

\noindent \textit{\textquotedblleft{}Você acha que é só uma história.\textquotedblright{}}

\noindent Você negou. Ou tentou. Ou apenas pensou em negar.

\noindent \textit{\textquotedblleft{}Você já sonhou comigo antes de ler o livro.\textquotedblright{}}

\noindent Era verdade. Você sonhou com olhos amarelos dias antes de encontrar este livro. Você tinha esquecido. O sonho trouxe de volta.

\noindent \textit{\textquotedblleft{}Eu estava esperando você.\textquotedblright{}}

\noindent Você acordou.

\noindent O quarto estava escuro. O relógio marcava 3:14 da manhã. A mesma hora todas as noites. Você não sabia disso ainda, mas vai aprender.

\vspace{0.5em}\noindent\rule{\textwidth}{0.5pt}\vspace{0.5em}

\subsection*{Noite 2}

\noindent Você tentou não dormir.

\noindent Ficou sentado na cama, com as costas encostadas na parede, a luz acesa, o livro fechado sobre a mesa. Você jurou que não ia ler antes de dormir. Jurou que ia pensar em outra coisa. Que ia ver televisão. Que ia ouvir música. Qualquer coisa menos a história.

\noindent Não adiantou.

\noindent O sono veio como um apagão. Você estava acordado num segundo e no sonho no segundo seguinte.

\noindent Dessa vez você estava na vala.

\noindent Você sabia que era a vala porque o cheiro era insuportável. O cheiro doce. O cheiro de flor murcha. Estava escuro. Seus pés estavam descalços. A terra estava fria e úmida.

\noindent Você olhou para baixo.

\noindent Seus pés estavam enterrados até o tornozelo.

\noindent Não em terra.

\noindent Em \textbf{corpos}.

\noindent Braços. Pernas. Troncos. Cabeças. Os corpos estavam empilhados. Alguns estavam duros. Outros estavam moles demais. A carne cedia sob seu peso como gelatina morna.

\noindent Você tentou se mexer. Seu pé direito afundou. Algo se mexeu debaixo de você.

\noindent \textit{\textquotedblleft{}Não se mexe.\textquotedblright{}}

\noindent A voz vinha de todos os lugares. Debaixo dos corpos. De dentro deles. De dentro de você.

\noindent \textit{\textquotedblleft{}Fica quieto. Ela está vindo.\textquotedblright{}}

\noindent \textit{Ela.} A criatura.

\noindent A pilha de corpos à sua frente começou a se mover. Os corpos rolaram para os lados. Braços balançaram. Uma mão --- não, uma garra --- emergiu do meio da massa.

\noindent Você queria correr. Não conseguia. Seus pés estavam presos. Os corpos seguravam você. Mãos mortas apertavam seus tornozelos.

\noindent A criatura subiu.

\noindent Você viu os espinhos de osso. A gaiola de costelas abertas. O órgão pulsando dentro, exposto, vivo, \textit{olhando} para você.

\noindent A cabeça apareceu. A face quase humana. A boca cheia de dentes demais.

\noindent Ela se inclinou sobre você.

\noindent O hálito quente no seu rosto. O cheiro adocicado. A boca se abriu.

\noindent \textit{\textquotedblleft{}Você não está lendo a história.\textquotedblright{}}

\noindent Ela colocou a mão na sua cabeça. Os dedos longos. As articulações extras. A unha preta arranhando seu couro cabeludo.

\noindent \textit{\textquotedblleft{}A história está lendo você.\textquotedblright{}}

\vspace{0.5em}\noindent\rule{\textwidth}{0.5pt}\vspace{0.5em}

\subsection*{Noite 3 --- Hoje}

\noindent Você está lendo este fragmento.

\noindent Já é noite. Ou está ficando. Você está lendo com pouca luz. O livro pesa nas suas mãos. O cheiro está aqui --- você já se acostumou, mas ele está aqui.

\noindent Você sabe que vai dormir esta noite. Sabe que vai sonhar. Sabe o que vai ver.

\noindent O olho.

\noindent A vala.

\noindent A criatura.

\noindent Você sabe que vai acordar às 3:14 da manhã. Vai olhar para o canto do quarto. E, desta vez, pode ser que não esteja vazio.

\noindent Você está com medo.

\noindent Bom.

\noindent O medo é o primeiro critério.

\vspace{0.5em}\noindent\rule{\textwidth}{0.5pt}\vspace{0.5em}

\noindent Você não vai fazer nada que este fragmento pede.

\noindent Acha que é sugestão. Acha que não funciona com você.

\noindent Funciona.

\noindent Esta noite, quando estiver sozinho, com o livro aberto na mesa de cabeceira e a luz apagada, você vai lembrar desta página. E vai olhar para o canto do quarto.

\noindent Vai ver o brilho.

\noindent E vai lembrar que a criatura nunca dorme. Ela só espera.

\vspace{0.5em}\noindent\rule{\textwidth}{0.5pt}\vspace{0.5em}

\subsection*{Noite 4}

\noindent Você não sonhou.

\noindent Nada aconteceu.

\noindent Você acordou descansado. O quarto estava normal. O cheiro tinha sumido. Você tomou café, saiu, viveu o dia inteiro sem pensar no livro.

\noindent À noite, deitou e dormiu como um menino.

\noindent Você achou que tinha vencido.

\noindent No fim da tarde do dia seguinte, sentou na sala e percebeu: a ausência do medo era pior que o medo. O silêncio era pior que o som. Você passou o dia inteiro esperando o susto que não veio. E, por não ter vindo, você passou o dia inteiro com medo do susto que poderia vir.

\noindent A criatura não tinha ido embora.

\noindent Ela tinha aprendido a se esconder melhor.

\noindent \textit{É mais aterrorizante quando ela não está --- porque você passa o tempo todo esperando.}

\noindent \textit{A espera é o verdadeiro contágio.}

\vspace{0.5em}\noindent\rule{\textwidth}{0.5pt}\vspace{0.5em}

\noindent \textit{--- O diagnóstico está preenchido. Falta apenas a assinatura.}

\noindent \textit{Passe o dedo na página. Na linha abaixo. Como se fosse assinar.}

\noindent \textit{--- Obrigado, paciente 1.263. O ciclo continua.}

\vspace{0.5em}
\noindent\textit{\textbf{↪ Rotas:} \rota{CONT-05} \rotasep \rota{MEM-15}}